\section{逆射与隐射定理的另一精美证法}\label{sec:8.5}

上节应用数学归纳法先证隐射定理: 方程组 $F(x,y) = 0$ 在条件 $\det J_{y}F \neq 0$ 时, 局部解得 $y = f(x)$, 使得 $F(x,f(x)) \equiv 0$. 进而, 当 $\det Jf(x) \neq 0$ 时, 对 $F(x,y) = f(x) - y = 0$ 施用隐射定理, 证得 $y = f(x)$ 有局部逆射 $x = g(y) = f^{-1}(y)$. 读者已领略了它的精美的论证方法. 自然会大胆地问: 能否先证局部逆射定理, 再应用局部逆射定理来证明隐射定理? 现在来实现这一设想.

\begin{theorem}[局部逆射定理]\label{thm:8.5.1}
设 $D \subset \mathbb{R}^{n}$ 为开集, $f: D \to \mathbb{R}^{n}$ 为 $C^{k}$ $(1 \leqslant k \leqslant +\infty)$ 映射, 且
\[
  \det Jf(\boldsymbol{x}^{0}) = \det\left(\frac{\partial f_{i}}{\partial x_{j}}(\boldsymbol{x}^{0})\right) \neq 0, \quad \boldsymbol{x}^{0} \in D,
\]
则 $f$ 为一个 $C^{k}$ 局部微分同胚, 它将 $\boldsymbol{x}^{0}$ 的一个开邻域映成 $f(\boldsymbol{x}^{0})$ 的一个开邻域.
\end{theorem}

\begin{proof}
可以假定 $\boldsymbol{x}^{0} = \boldsymbol{0}$, $f(\boldsymbol{x}^{0}) = \boldsymbol{0}$, $\left(\dfrac{\partial f_{i}}{\partial x_{j}}(\boldsymbol{x}^{0})\right) = (\delta_{j}^{i}) = \boldsymbol{I}$ (否则作一仿射变换, 不影响定理的证明). 选 $a$ 充分小, 使得方体
\[
  C^{n}(a) = \bigl\{\boldsymbol{x} \in \mathbb{R}^{n} \mid \|\boldsymbol{x}\| = \max_{1 \leqslant j \leqslant n}|x_{j}| < a\bigr\} \subset D,
\]
且当 $\boldsymbol{x} \in C^{n}(a)$ 时,
\[
  \det\left(\frac{\partial f_{i}}{\partial x_{j}}(\boldsymbol{x})\right) \neq 0
\]
及
\[
  \left(\frac{\partial f_{i}}{\partial x_{j}}(\boldsymbol{x})\right) - \boldsymbol{I}
\]
中的元素满足
\[
  \left|\left(\frac{\partial f_{i}}{\partial x_{j}}(\boldsymbol{x})\right) - \delta_{j}^{i}\right| \leqslant \frac{1}{2n}
  \quad \biggl(\text{只需注意}\ \left(\frac{\partial f_{i}}{\partial x_{j}}(\boldsymbol{0})\right) - \boldsymbol{I} = \boldsymbol{0}\biggr).
\]
令 $g(\boldsymbol{x}) = f(\boldsymbol{x}) - \boldsymbol{x}$, 则当 $\boldsymbol{x}^{1}, \boldsymbol{x}^{2} \in C^{n}(a)$ 时, 应用 Lagrange 中值定理有
\begin{align*}
  \|\boldsymbol{g}(\boldsymbol{x}^{1}) - \boldsymbol{g}(\boldsymbol{x}^{2})\| &= \max_{1 \leqslant i \leqslant n}\bigl\{|g_{i}(\boldsymbol{x}^{1}) - g_{i}(\boldsymbol{x}^{2})|\bigr\} \\
  &= \max_{1 \leqslant i \leqslant n}\left\{\left|\sum_{j=1}^{n}\frac{\partial g_{i}}{\partial x_{j}}(\boldsymbol{\xi}^{j})(x_{j}^{1} - x_{j}^{2})\right|\right\} \\
  &\leqslant n\max_{1 \leqslant i, j \leqslant n}\left\{\left|\frac{\partial g_{i}}{\partial x_{j}}(\boldsymbol{\xi}^{j})\right|\right\}\max_{1 \leqslant j \leqslant n}\bigl\{|x_{j}^{1} - x_{j}^{2}|\bigr\} \\
  &\leqslant n \cdot \frac{1}{2n}\|\boldsymbol{x}^{1} - \boldsymbol{x}^{2}\| = \frac{1}{2}\|\boldsymbol{x}^{1} - \boldsymbol{x}^{2}\|.
\end{align*}
此外, 还有
\begin{align*}
  \|f(\boldsymbol{x}^{1}) - f(\boldsymbol{x}^{2})\| &= \|\boldsymbol{g}(\boldsymbol{x}^{1}) - \boldsymbol{g}(\boldsymbol{x}^{1}) + \boldsymbol{x}^{1} - \boldsymbol{x}^{2}\| \\
  &\geqslant \|\boldsymbol{x}^{1} - \boldsymbol{x}^{2}\| - \|\boldsymbol{g}(\boldsymbol{x}^{1}) - \boldsymbol{g}(\boldsymbol{x}^{2})\| \\
  &\geqslant \|\boldsymbol{x}^{1} - \boldsymbol{x}^{2}\| - \frac{1}{2}\|\boldsymbol{x}^{1} - \boldsymbol{x}^{2}\| \\
  &= \frac{1}{2}\|\boldsymbol{x}^{1} - \boldsymbol{x}^{2}\|.
\end{align*}
\begin{enumerate}
\item[(1)] 如果 $\boldsymbol{y} \in C^{n}\left(\dfrac{a}{2}\right)$, 则必存在一点 $\boldsymbol{x} \in C^{n}(a)$, 使得 $f(\boldsymbol{x}) = \boldsymbol{y}$.

为此, 设
\[
  \boldsymbol{x}^{0} = \boldsymbol{0}, \quad \boldsymbol{x}^{1} = \boldsymbol{y}, \quad \boldsymbol{x}^{k+1} = \boldsymbol{y} - \boldsymbol{g}(\boldsymbol{x}^{k}), \quad k = 0, 1, 2, \cdots,
\]
容易看出
\begin{align*}
  \|\boldsymbol{x}^{k} - \boldsymbol{x}^{k-1}\| &= \|\boldsymbol{g}(\boldsymbol{x}^{k-1}) - \boldsymbol{g}(\boldsymbol{x}^{k-2})\| \leqslant \frac{1}{2}\|\boldsymbol{x}^{k-1} - \boldsymbol{x}^{k-2}\| \\
  \leqslant {} & \cdots \leqslant \frac{1}{2^{k-2}}\|\boldsymbol{x}^{1} - \boldsymbol{x}^{0}\| = \frac{1}{2^{k-1}}\|\boldsymbol{y}\|,
\end{align*}
\begin{align*}
  \|\boldsymbol{x}^{k}\| &= \|\boldsymbol{x}^{k} - \boldsymbol{x}^{0}\| \\
  &\leqslant \|\boldsymbol{x}^{k} - \boldsymbol{x}^{k-1}\| + \|\boldsymbol{x}^{k-1} - \boldsymbol{x}^{k-2}\| + \cdots + \|\boldsymbol{x}^{1} - \boldsymbol{x}^{0}\| \\
  &\leqslant \left(\frac{1}{2^{k-1}} + \frac{1}{2^{k-2}} + \cdots + 1\right)\|\boldsymbol{y}\| < 2\|\boldsymbol{y}\| < 2 \cdot \frac{a}{2} = a,
\end{align*}
\begin{align*}
  \|\boldsymbol{x}^{k+p} - \boldsymbol{x}^{k}\| &\leqslant \|\boldsymbol{x}^{k+p} - \boldsymbol{x}^{k+p-1}\| + \cdots + \|\boldsymbol{x}^{k+1} - \boldsymbol{x}^{k}\| \\
  &\leqslant \left(\frac{1}{2^{k+p-1}} + \cdots + \frac{1}{2^{k}}\right)\|\boldsymbol{y}\| \\
  &= \frac{1}{2^{k}}\left(\frac{1}{2^{p-1}} + \cdots + 1\right)\|\boldsymbol{y}\| \leqslant \frac{1}{2^{k-1}}\|\boldsymbol{y}\|.
\end{align*}
于是, $\boldsymbol{x}^{k} \in C^{n}(a)$ 及 $\{\boldsymbol{x}^{k}\}$ 为 Cauchy 点列. 由于 $\mathbb{R}^{n}$ 完备, 故 $\{\boldsymbol{x}^{k}\}$ 收敛, 设 $\lim\limits_{k \to +\infty}\boldsymbol{x}^{k} = \boldsymbol{x}$, 再在上述第二个不等式中令 $k \to +\infty$ 得到
\[
  \|\boldsymbol{x}\| < 2\|\boldsymbol{y}\| < a,
\]
因而 $\boldsymbol{x} \in C^{n}(a)$, 且
\begin{align*}
  \boldsymbol{x} &= \lim_{k \to +\infty}\boldsymbol{x}^{k+1} = \lim_{k \to +\infty}\bigl[\boldsymbol{y} - \boldsymbol{g}(\boldsymbol{x}^{k})\bigr] = \boldsymbol{y} - \boldsymbol{g}(\boldsymbol{x}) \\
  &= \boldsymbol{y} - \bigl[f(\boldsymbol{x}) - \boldsymbol{x}\bigr] = \boldsymbol{y} - f(\boldsymbol{x}) + \boldsymbol{x}, \quad f(\boldsymbol{x}) = \boldsymbol{y}.
\end{align*}
这就证明了 $\boldsymbol{x}$ 的存在性.

再证惟一性. 如果 $\exists\,\boldsymbol{x}, \boldsymbol{x}' \in C^{n}(a)$, s. t.\ $f(\boldsymbol{x}) = \boldsymbol{y} = f(\boldsymbol{x}')$, 则由
\[
  0 = \|f(\boldsymbol{x}) - f(\boldsymbol{x}')\| \geqslant \frac{1}{2}\|\boldsymbol{x} - \boldsymbol{x}'\| \geqslant 0,
\]
推得 $\|\boldsymbol{x} - \boldsymbol{x}'\| = 0$, $\boldsymbol{x} = \boldsymbol{x}'$.

\item[(2)] 由 (1) 知 $f^{-1}: C^{n}\left(\dfrac{a}{2}\right) \to C^{n}(a)$ 是存在的. 另一方面, 从
\[
  \|\boldsymbol{y}' - \boldsymbol{y}''\| = \|f(\boldsymbol{x}') - f(\boldsymbol{x}'')\| \geqslant \frac{1}{2}\|\boldsymbol{x}' - \boldsymbol{x}''\| = \frac{1}{2}\|f^{-1}(\boldsymbol{y}') - f^{-1}(\boldsymbol{y}'')\|
\]
可推出 $f^{-1}$ 在 $C^{n}\left(\dfrac{a}{2}\right)$ 上是一致连续的, 因而是连续的.

显然, 由 $f$ 连续及 $C^{n}\left(\dfrac{a}{2}\right)$ 为开集知, $f^{-1}\left(C^{n}\left(\dfrac{a}{2}\right)\right)$ 为开集, 从而 $C^{n}\left(\dfrac{a}{2}\right)$ 在 $f^{-1}: C^{n}\left(\dfrac{a}{2}\right) \to C^{n}(a)$ 下的逆像为 $C^{n}(a) \cap f^{-1}\left(C^{n}\left(\dfrac{a}{2}\right)\right)$, 它是开集. 这就证明了
\[
  f: C^{n}(a) \cap f^{-1}\left(C^{n}\left(\frac{a}{2}\right)\right) \to C^{n}\left(\frac{a}{2}\right)
\]
为一个同胚 (拓扑映射), 它将 $\boldsymbol{0}$ 的一个开邻域映成 $f(\boldsymbol{0}) = \boldsymbol{0}$ 的一个开邻域.

\item[(3)] $f^{-1}$ 是可微的, 由此得到 $J(f^{-1})$ 是存在的.

由于 $f$ 是 $C^{k}$ $(1 \leqslant k \leqslant +\infty)$, 根据\cref{thm:8.2.4}, $f$ 在 $\forall\,\boldsymbol{x}^{1} \in C^{n}(a) \cap f^{-1}\left(C^{n}\left(\dfrac{a}{2}\right)\right)$ 是可微的, 即
\[
  f(\boldsymbol{x}) - f(\boldsymbol{x}') = \left(\frac{\partial f_{i}}{\partial x_{j}}(\boldsymbol{x}^{1})\right)(\boldsymbol{x} - \boldsymbol{x}^{1}) + \boldsymbol{h}(\boldsymbol{x}^{1}, \boldsymbol{x}),
\]
\[
  \lim_{\|\boldsymbol{x} - \boldsymbol{x}^{1}\| \to 0}\frac{\boldsymbol{h}(\boldsymbol{x}^{1}, \boldsymbol{x})}{\|\boldsymbol{x} - \boldsymbol{x}^{1}\|} = 0
  \quad (\boldsymbol{x}^{1}\ \text{固定}), \quad \boldsymbol{x}^{1}, \boldsymbol{x} \in C^{n}(a) \cap \boldsymbol{f}^{-1}\left(C^{n}\left(\frac{a}{2}\right)\right).
\]
记 $\left(\dfrac{\partial f_{i}}{\partial x_{j}}(\boldsymbol{x}^{1})\right)^{-1}$ 为 $\left(\dfrac{\partial f_{i}}{\partial x_{j}}(\boldsymbol{x}^{1})\right)$ 的逆矩阵, 则
\[
  \left(\frac{\partial f_{i}}{\partial x_{j}}(\boldsymbol{x}^{1})\right)^{-1}\bigl[\boldsymbol{f}(\boldsymbol{x}) - \boldsymbol{f}(\boldsymbol{x}^{1})\bigr] = (\boldsymbol{x} - \boldsymbol{x}^{1}) + \left(\frac{\partial f_{i}}{\partial x_{j}}(\boldsymbol{x}^{1})\right)^{-1}\boldsymbol{h}(\boldsymbol{x}^{1}, \boldsymbol{x}),
\]
\[
  \left(\frac{\partial f_{i}}{\partial x_{j}}(\boldsymbol{x}^{1})\right)^{-1}(\boldsymbol{y} - \boldsymbol{y}^{1}) + \tilde{\boldsymbol{h}}(\boldsymbol{y}^{1}, \boldsymbol{y}) = \boldsymbol{f}^{-1}(\boldsymbol{y}) - \boldsymbol{f}^{-1}(\boldsymbol{y}^{1}),
\]
其中 $\tilde{\boldsymbol{h}}(\boldsymbol{y}^{1}, \boldsymbol{y}) = -\left(\dfrac{\partial f_{i}}{\partial x_{j}}(\boldsymbol{x}^{1})\right)^{-1}\boldsymbol{h}(\boldsymbol{x}^{1}, \boldsymbol{x})$. 因为
\[
  \|\boldsymbol{y} - \boldsymbol{y}^{1}\| = \|f(\boldsymbol{x}) - f(\boldsymbol{x}^{1})\| \geqslant \frac{1}{2}\|\boldsymbol{x} - \boldsymbol{x}^{1}\|,
\]
\[
  \frac{\|\boldsymbol{x} - \boldsymbol{x}^{1}\|}{\|\boldsymbol{y} - \boldsymbol{y}^{1}\|} \leqslant 2, \quad \|\boldsymbol{y} - \boldsymbol{y}^{1}\| \neq 0,
\]
所以,
\begin{align*}
  \lim_{\boldsymbol{y} \to \boldsymbol{y}^{1}}\frac{\tilde{\boldsymbol{h}}(\boldsymbol{y}^{1}, \boldsymbol{y})}{\|\boldsymbol{y} - \boldsymbol{y}^{1}\|} &= \lim_{\boldsymbol{y} \to \boldsymbol{y}^{1}}\frac{-\left(\frac{\partial f_{i}}{\partial x_{j}}(\boldsymbol{x}^{1})\right)^{-1}\boldsymbol{h}(f^{-1}(\boldsymbol{y}^{1}), f^{-1}(\boldsymbol{y}))}{\|\boldsymbol{y} - \boldsymbol{y}^{1}\|} \\
  &= -\left(\frac{\partial f_{i}}{\partial x_{j}}(\boldsymbol{x}^{1})\right)^{-1}\lim_{\boldsymbol{x} \to \boldsymbol{x}^{1}}\frac{\boldsymbol{h}(\boldsymbol{x}^{1}, \boldsymbol{x})}{\|\boldsymbol{x} - \boldsymbol{x}^{1}\|} \cdot \frac{\|\boldsymbol{x} - \boldsymbol{x}^{1}\|}{\|\boldsymbol{y} - \boldsymbol{y}^{1}\|} = \boldsymbol{0}.
\end{align*}
{\emergencystretch=3em 这就证明了 $f^{-1}$ 在 $\boldsymbol{y}^{1}$ 是可微的. 根据\cref{thm:8.2.3} 可推出 $(f^{-1})_{i}$ 在 $\boldsymbol{y}^{1}$ 有关于 $y_{1}, y_{2}, \cdots, y_{n}$ 的偏导数, 且\par}
\[
  \left(\frac{\partial (f^{-1})_{i}}{\partial y_{j}}(\boldsymbol{y}^{1})\right) = \left(\frac{\partial f_{i}}{\partial x_{j}}(\boldsymbol{x}^{1})\right)^{-1}.
\]

\item[(4)] 最后, 用数学归纳法证明 $f^{-1}$ 是 $C^{k}$ 的. 映射
\[
  \boldsymbol{y} \mapsto J(f^{-1})(\boldsymbol{y}) = \left(\frac{\partial (f^{-1})_{i}}{\partial y_{j}}(\boldsymbol{y})\right) = \left(\frac{\partial f_{i}}{\partial x_{j}}(\boldsymbol{x})\right)^{-1} = (Jf(\boldsymbol{x}))^{-1}
\]
可以由以下映射的复合而得到:
\begin{center}
\begin{tikzpicture}[x=0.98cm, y=0.80cm, >=latex, line cap=round, line join=round]
  \node[inner sep=1pt] (A) at (0, 0) {$C^{n}\left(\dfrac{a}{2}\right)$};
  \node[inner sep=1pt] (B) at (2.60, 0) {$C^{n}(a)$};
  \node[inner sep=1pt] (C) at (5.05, 0) {$GL(n, \mathbb{R})$};
  \node[inner sep=1pt] (D) at (8.45, 0) {$GL(n, \mathbb{R})$};
  \draw[->, line width=0.7pt] (A.east) -- node[above=1.5pt, inner sep=1pt] {$f^{-1}$} (B.west);
  \draw[->, line width=0.7pt] (B.east) -- node[above=1.5pt, inner sep=1pt] {$Jf$} (C.west);
  \draw[->, line width=0.7pt] (C.east) -- node[above=1.5pt, inner sep=1pt] {矩阵的逆} (D.west);
  \draw[decorate, decoration={brace, amplitude=4.5pt, mirror}] (A.south) -- (D.south);
  \node[below=6.5pt, inner sep=1pt] at (4.22, -0.42) {$J(f^{-1})$};
  \node[right=2pt, inner sep=1pt] at (D.east) {$,$};
\end{tikzpicture}
\end{center}
其中 $GL(n, \mathbb{R})$ 表示 $n \times n$ 的非异矩阵的全体. 根据 (2), $f^{-1}$ 是 $C^{0}$ 的 (即是连续的). 此外, 显然 $Jf$ 是 $C^{k-1}$ 的, 矩阵的逆所对应的映射是 $C^{\infty}$ 的, 所以 $J(f^{-1})$ 是 $C^{0}$ 的, 即 $f^{-1}$ 是 $C^{1}$ 的. 如果 $f^{-1}$ 是 $C^{s}$ 的 $(1 \leqslant s \leqslant k-1)$, 则同样的推理可得到 $J(f^{-1})$ 是 $C^{s}$ 的, 即 $f^{-1}$ 是 $C^{s+1}$ 的. 这就证明了 $f^{-1}$ 是 $C^{k}$ 的.
\end{enumerate}
\end{proof}
应用逆射定理可以证明隐射定理.

\begin{theorem}[隐射定理]\label{thm:8.5.2}
设 $(\boldsymbol{x}^{0}, \boldsymbol{y}^{0}) \in U \subset \mathbb{R}^{n} \times \mathbb{R}^{m} = \mathbb{R}^{n+m}$, $U$ 为开集. $F: U \to \mathbb{R}^{m}$ 为 $C^{k}$ $(1 \leqslant k \leqslant +\infty)$ 映射, $F(\boldsymbol{x}^{0}, \boldsymbol{y}^{0}) = \boldsymbol{0}$, 且
\[
  \left.\frac{\partial(F_{1}, F_{2}, \cdots, F_{m})}{\partial(y_{1}, y_{2}, \cdots, y_{m})}\right|_{(\boldsymbol{x}^{0}, \boldsymbol{y}^{0})} \neq 0,
\]
则存在惟一的 $C^{k}$ 映射 $f: I \to \mathbb{R}^{m}$ (其中 $I$ 为 $\boldsymbol{x}^{0}$ 的一个开邻域), 使得 $\boldsymbol{y} = f(\boldsymbol{x})$, $\boldsymbol{y}^{0} = f(\boldsymbol{x}^{0})$, 且
\[
  \boldsymbol{F}(\boldsymbol{x}, f(\boldsymbol{x})) = \boldsymbol{0}, \quad \boldsymbol{x} \in I.
\]
\end{theorem}

\begin{proof}
定义 $\widetilde{F}: U \to \mathbb{R}^{n+m}$ 为
\[
  \widetilde{\boldsymbol{F}}(\boldsymbol{x}, \boldsymbol{y}) = \bigl(\boldsymbol{x}, \boldsymbol{F}(\boldsymbol{x}, \boldsymbol{y})\bigr).
\]
显然,
\begin{align*}
  \det\bigl(J\widetilde{\boldsymbol{F}}(\boldsymbol{x}, \boldsymbol{y})\bigr)_{(\boldsymbol{x}^{0}, \boldsymbol{y}^{0})} &= \det\begin{pmatrix}
    \boldsymbol{I}_{n} & \boldsymbol{0} \\
    \boldsymbol{J}_{\boldsymbol{x}}\boldsymbol{F} & \boldsymbol{J}_{\boldsymbol{y}}\boldsymbol{F}
  \end{pmatrix}_{(\boldsymbol{x}^{0}, \boldsymbol{y}^{0})} \\
  &= \left.\frac{\partial(F_{1}, F_{2}, \cdots, F_{m})}{\partial(y_{1}, y_{2}, \cdots, y_{m})}\right|_{(\boldsymbol{x}^{0}, \boldsymbol{y}^{0})} \neq 0,
\end{align*}
因此, 在 $(\boldsymbol{x}^{0}, \boldsymbol{y}^{0})$ 邻近有局部逆射 $\widetilde{G}$, 它是 $C^{k}$ 映射.

如果 $F(\boldsymbol{x}, f(\boldsymbol{x})) = \boldsymbol{0}$, 则必须有
\[
  \widetilde{\boldsymbol{F}}(\boldsymbol{x}, f(\boldsymbol{x})) = \bigl(\boldsymbol{x}, F(\boldsymbol{x}, f(\boldsymbol{x}))\bigr) = (\boldsymbol{x}, \boldsymbol{0}),
\]
\[
  \bigl(\boldsymbol{x}, f(\boldsymbol{x})\bigr) = \widetilde{G} \circ \widetilde{F}(\boldsymbol{x}, f(\boldsymbol{x})) = \widetilde{G}(\boldsymbol{x}, \boldsymbol{0}),
\]
\[
  f(\boldsymbol{x}) = \pi \circ \widetilde{G}(\boldsymbol{x}, \boldsymbol{0}),
\]
其中 $\pi: \mathbb{R}^{n} \times \mathbb{R}^{m} \to \mathbb{R}^{m}$, $\pi(\boldsymbol{x}, \boldsymbol{y}) = \boldsymbol{y}$.

另一方面, 若令 $f(\boldsymbol{x}) = \pi \circ \widetilde{G}(\boldsymbol{x}, \boldsymbol{0})$, 显然它是 $C^{k}$ 的, 且
\[
  \bigl(\boldsymbol{x}, \pi \circ \widetilde{G}(\boldsymbol{x}, \boldsymbol{0})\bigr) = \bigl(\boldsymbol{x}, f(\boldsymbol{x})\bigr) = \widetilde{G}(\boldsymbol{x}, \boldsymbol{0}),
\]
\begin{align*}
  \bigl(\boldsymbol{x}, F(\boldsymbol{x}, f(\boldsymbol{x}))\bigr) &= \widetilde{\boldsymbol{F}}(\boldsymbol{x}, f(\boldsymbol{x})) = \widetilde{\boldsymbol{F}}\bigl(\widetilde{G}(\boldsymbol{x}, \boldsymbol{0})\bigr) \\
  &= \widetilde{\boldsymbol{F}} \circ \widetilde{G}(\boldsymbol{x}, \boldsymbol{0}) = (\boldsymbol{x}, \boldsymbol{0}),
\end{align*}
\[
  \boldsymbol{F}(\boldsymbol{x}, f(\boldsymbol{x})) = \boldsymbol{0}.
\]
再由
\begin{align*}
  \widetilde{F}(\boldsymbol{x}^{0}, \boldsymbol{y}^{0}) &= (\boldsymbol{x}^{0}, F(\boldsymbol{x}^{0}, \boldsymbol{y}^{0})) = (\boldsymbol{x}^{0}, \boldsymbol{0}) \\
  &= \bigl(\boldsymbol{x}^{0}, F(\boldsymbol{x}^{0}, f(\boldsymbol{x}^{0}))\bigr) = \widetilde{F}(\boldsymbol{x}^{0}, f(\boldsymbol{x}^{0}))
\end{align*}
推出
\[
  \bigl(\boldsymbol{x}^{0}, \boldsymbol{y}^{0}\bigr) = \widetilde{\boldsymbol{G}} \circ \widetilde{\boldsymbol{F}}(\boldsymbol{x}^{0}, \boldsymbol{y}^{0}) = \widetilde{\boldsymbol{G}} \circ \widetilde{\boldsymbol{F}}\bigl(\boldsymbol{x}^{0}, f(\boldsymbol{x}^{0})\bigr) = \bigl(\boldsymbol{x}^{0}, f(x_{0})\bigr),
\]
\[
  \boldsymbol{y}^{0} = f(\boldsymbol{x}^{0}),
\]
所以, $f(\boldsymbol{x}) = \pi \cdot \widetilde{G}(\boldsymbol{x}, \boldsymbol{0})$ 就是所要求的惟一的 $C^{k}$ 映射.
\end{proof}

作为局部逆射定理的应用, 我们给出下面三个推论.

\begin{corollary}\label{cor:8.5.1}
设 $U$ 为 $\mathbb{R}^{n}$ 中的开集, $f: U \to \mathbb{R}^{m}$ 为 $C^{k}$ 映射 $(1 \leqslant k \leqslant +\infty)$, $f(\boldsymbol{0}) = \boldsymbol{0}$, 且
\[
  \operatorname{rank}\left(\frac{\partial f_{i}}{\partial x_{j}}(\boldsymbol{0})\right) = n \quad (n \leqslant m),
\]
则存在一个 $C^{k}$ 微分同胚 $g$, 将 $\mathbb{R}^{m}$ 中 $\boldsymbol{0}$ 的一个开邻域映成另一个开邻域, 且 $g(\boldsymbol{0}) = \boldsymbol{0}$ 及
\[
  \boldsymbol{g} \circ f(x_{1}, x_{2}, \cdots, x_{n}) = (x_{1}, x_{2}, \cdots, x_{n}, 0, 0, \cdots, 0).
\]
\end{corollary}

\begin{proof}
不失一般性, 可以假定
\[
  \left.\frac{\partial(f_{1}, f_{2}, \cdots, f_{n})}{\partial(x_{1}, x_{2}, \cdots, x_{n})}\right|_{\boldsymbol{0}} \neq 0.
\]
定义 $F: U \times \mathbb{R}^{m-n} \to \mathbb{R}^{m}$ 为
\[
  \boldsymbol{F}(x_{1}, x_{2}, \cdots, x_{m}) = f(x_{1}, x_{2}, \cdots, x_{n}) + (0, 0, \cdots, 0, x_{n+1} + x_{n+2} + \cdots, x_{m}),
\]
则 $F(\boldsymbol{0}) = \boldsymbol{0}$, 且
\[
  \left.\frac{\partial(f_{1}, f_{2}, \cdots, f_{m})}{\partial(x_{1}, x_{2}, \cdots, x_{m})}\right|_{\boldsymbol{0}} = \det\begin{pmatrix}
    Jf(\boldsymbol{0}) & \boldsymbol{0} \\
    & \boldsymbol{I}_{m-n}
  \end{pmatrix} = \left.\frac{\partial(f_{1}, f_{2}, \cdots, f_{n})}{\partial(x_{1}, x_{2}, \cdots, x_{n})}\right|_{\boldsymbol{0}} \neq 0,
\]
根据\cref{thm:8.5.1}, 存在 $F$ 的局部逆映射 $g$, 它是 $C^{k}$ 微分同胚, 将 $\mathbb{R}^{m}$ 中 $\boldsymbol{0}$ 的一个开邻域映成另一个开邻域, 使得 $g(\boldsymbol{0}) = \boldsymbol{0}$ 及
\[
  \boldsymbol{g} \circ f(x_{1}, x_{2}, \cdots, x_{n}) = \boldsymbol{g}\bigl(\boldsymbol{F}(x_{1}, x_{2}, \cdots, x_{n}, 0, 0, \cdots, 0)\bigr) = (x_{1}, x_{2}, \cdots, x_{n}, 0, 0, \cdots, 0).
\]
\end{proof}

\begin{corollary}\label{cor:8.5.2}
设 $U$ 为 $\mathbb{R}^{n}$ 中的开集, $f: U \to \mathbb{R}^{m}$ 为 $C^{k}$ 映射 $(1 \leqslant k \leqslant +\infty)$, $f(\boldsymbol{0}) = \boldsymbol{0}$, 且
\[
  \operatorname{rank}\left(\frac{\partial f_{i}}{\partial y_{j}}(\boldsymbol{0})\right) = m \quad (m \leqslant n),
\]
则存在一个 $C^{k}$ 微分同胚 $h$, 将 $\mathbb{R}^{n}$ 中 $\boldsymbol{0}$ 的一个开邻域映成另一个开邻域, 使得 $h(\boldsymbol{0}) = \boldsymbol{0}$ 及
\[
  \boldsymbol{f} \circ \boldsymbol{h}(x_{1}, x_{2}, \cdots, x_{n}) = (x_{1}, x_{2}, \cdots, x_{m}).
\]
\end{corollary}

\begin{proof}
因为 $\operatorname{rank}\left(\dfrac{\partial f_{i}}{\partial y_{j}}(\boldsymbol{0})\right) = m$, 不失一般性, 可以假定
\[
  \left.\frac{\partial(f_{1}, f_{2}, \cdots, f_{m})}{\partial(y_{1}, y_{2}, \cdots, y_{m})}\right|_{\boldsymbol{0}} \neq 0.
\]
定义 $F: U \to \mathbb{R}^{n}$ 为
\[
  \boldsymbol{F}(y_{1}, y_{2}, \cdots, y_{n}) = \bigl(f_{1}(\boldsymbol{y}), f_{2}(\boldsymbol{y}), \cdots, f_{m}(\boldsymbol{y}), y_{m+1}, y_{m+2}, \cdots, y_{n}\bigr),
\]
则 $\boldsymbol{F}(\boldsymbol{0}) = \boldsymbol{0}$, $f = \pi^{\circ} \boldsymbol{F}$ (其中 $\pi: \mathbb{R}^{n} \to \mathbb{R}^{m}$, $\pi(x_{1}, x_{2}, \cdots, x_{n}) = (x_{1}, x_{2}, \cdots, x_{m})$ 为投影), 且
\[
  \left.\frac{\partial(F_{1}, F_{2}, \cdots, F_{n})}{\partial(y_{1}, y_{2}, \cdots, y_{n})}\right|_{\boldsymbol{0}} = \det\begin{pmatrix}
    \boldsymbol{J}_{y} & f(\boldsymbol{0}) \\
    \boldsymbol{0} & \boldsymbol{I}_{n-m}
  \end{pmatrix} = \left.\frac{\partial(f_{1}, f_{2}, \cdots, f_{m})}{\partial(y_{1}, y_{2}, \cdots, y_{m})}\right|_{\boldsymbol{0}} \neq 0,
\]
根据\cref{thm:8.5.1}, 存在 $F$ 的局部逆映射 $h$, 它是 $C^{k}$ 微分同胚, 将 $\mathbb{R}^{n}$ 中 $\boldsymbol{0}$ 的一个开邻域映成另一个开邻域, 使得 $h(\boldsymbol{0}) = \boldsymbol{0}$ 及
\[
  f \circ h(x_{1}, x_{2}, \cdots, x_{h}) = \pi \circ F \circ h(x_{1}, x_{2}, \cdots, x_{n}) = \pi(x_{1}, x_{2}, \cdots, x_{n}) = (x_{1}, x_{2}, \cdots, x_{m}).
\]
\end{proof}

\begin{corollary}\label{cor:8.5.3}
设 $U$ 为 $\mathbb{R}^{n}$ 中的开集, $f: U \to \mathbb{R}^{m}$ 为 $C^{k}$ 映射 $(1 \leqslant k \leqslant +\infty)$, $f(\boldsymbol{0}) = \boldsymbol{0}$, $\operatorname{rank}\left(\dfrac{\partial f_{i}}{\partial y_{j}}(\boldsymbol{y})\right) = l$ $(l \leqslant \min\{m, n\})$, 则存在 $C^{k}$ 微分同胚 $h$ 与 $g$, 它们分别将 $\mathbb{R}^{n}$ 与 $\mathbb{R}^{m}$ 的 $\boldsymbol{0}$ 的一个开邻域映成 $\boldsymbol{0}$ 的另一个开邻域, 使得 $h(\boldsymbol{0}) = \boldsymbol{0}$, $g(\boldsymbol{0}) = \boldsymbol{0}$ 及
\[
  \boldsymbol{g} \circ \boldsymbol{f} \circ \boldsymbol{h}(x_{1}, x_{2}, \cdots, x_{n}) = (x_{1}, x_{2}, \cdots, x_{l}, 0, 0, \cdots, 0).
\]
\end{corollary}

\begin{proof}
不失一般性, 不妨设
\[
  \left.\frac{\partial(f_{1}, f_{2}, \cdots, f_{l})}{\partial(y_{1}, y_{2}, \cdots, y_{l})}\right|_{\boldsymbol{0}} \neq 0.
\]
令
\begin{align*}
  (x_{1}, x_{2}, \cdots, x_{l}, \cdots, x_{n}) &= \boldsymbol{F}(y_{1}, y_{2}, \cdots, y_{n}) \\
  &= \bigl(f_{1}(\boldsymbol{y}), f_{2}(\boldsymbol{y}), \cdots, f_{l}(\boldsymbol{y}), y_{l+1}, y_{l+2}, \cdots, y_{n}\bigr),
\end{align*}
\[
  \boldsymbol{F}(\boldsymbol{0}) = \boldsymbol{0},
\]
\[
  \left.\frac{\partial(f_{1}, f_{2}, \cdots, f_{l})}{\partial(y_{1}, y_{2}, \cdots, y_{l})}\right|_{\boldsymbol{0}} = \det\begin{pmatrix}
    \dfrac{\partial f_{1}}{\partial y_{1}} & \cdots & \dfrac{\partial f_{1}}{\partial y_{l}} & & \\
    \vdots & & \vdots & & \\
    \dfrac{\partial f_{l}}{\partial y_{1}} & \cdots & \dfrac{\partial f_{l}}{\partial y_{l}} & & \\
    \boldsymbol{0} & & \boldsymbol{I}_{n-l} & &
  \end{pmatrix}_{\boldsymbol{0}} = \left.\frac{\partial(f_{1}, f_{2}, \cdots, f_{l})}{\partial(y_{1}, y_{2}, \cdots, y_{l})}\right|_{\boldsymbol{0}} \neq 0.
\]
由\cref{thm:8.5.1}, $F$ 有局部逆映射 $h$, 将 $\mathbb{R}^{n}$ 中 $\boldsymbol{0}$ 的一个凸开邻域 $U_{1} \subset U$ 映成 $\boldsymbol{0}$ 的另一个开邻域, 且
\begin{align*}
  (z_{1}, z_{2}, \cdots, z_{m}) &= f \circ h(\boldsymbol{x}) = \bigl(f_{1} \circ h(\boldsymbol{x}), f_{2} \circ h(\boldsymbol{x}), \cdots, f_{l} \circ h(\boldsymbol{x}), f_{l+1} \circ h(\boldsymbol{x}), \cdots, f_{m} \circ h(\boldsymbol{x})\bigr) \\
  &= \bigl(f_{1}(\boldsymbol{y}), f_{2}(\boldsymbol{y}), \cdots, f_{l}(\boldsymbol{y}), f_{l+1} \circ h(\boldsymbol{x}), \cdots, f_{m} \circ h(\boldsymbol{x})\bigr) \\
  &= \bigl(x_{1}, x_{2}, \cdots, x_{l}, f_{l+1} \circ h(\boldsymbol{x}), \cdots, f_{m} \circ h(\boldsymbol{x})\bigr).
\end{align*}
因为
\[
  \left.\operatorname{rank}\left(\frac{\partial (f \circ h)_{i}}{\partial x_{j}}\right)\right|_{U_{1}} \equiv l,
\]
所以
\[
  0 = \left.\frac{\partial(z_{1}, z_{2}, \cdots, z_{l}, \cdots, z_{l+j})}{\partial(x_{1}, x_{2}, \cdots, x_{l}, \cdots, x_{l+s})}\right|_{U_{1}} = \left.\det\begin{pmatrix}
    \boldsymbol{I}_{l} & \boldsymbol{0} \\
    * & \dfrac{\partial(f_{l+j} \circ h)}{\partial x_{l+s}}
  \end{pmatrix}\right|_{U_{1}}.
\]
于是, $z_{l+j} = f_{l+j} \circ h(\boldsymbol{x})$ 仅是 $x_{1}, x_{2}, \cdots, x_{l}$ 的函数, 记为
\[
  z_{l+j} = f_{l+j} \circ h(\boldsymbol{x}) = \widetilde{f}_{l+j}(x_{1}, x_{2}, \cdots, x_{l}), \quad j = 1, 2, \cdots, m - l.
\]
令
\begin{align*}
  &\bigl(z_{1}, z_{2}, \cdots, z_{l}, z_{l+1}, \cdots, z_{m}\bigr) \\
  ={} & \boldsymbol{G}(u_{1}, u_{2}, \cdots, u_{m}) \\
  ={} & \bigl(u_{1}, u_{2}, \cdots, u_{l}, \widetilde{f}_{l+1}(u_{1}, u_{2}, \cdots, u_{l}), \cdots, \widetilde{f}_{m}(u_{1}, u_{2}, \cdots, u_{l})\bigr) \\
  &\quad + (0, 0, \cdots, 0, u_{l+1}, u_{l+2}, \cdots, u_{m}),
\end{align*}
\[
  \boldsymbol{G}(\boldsymbol{0}) = \boldsymbol{0},
\]
\[
  \boldsymbol{G}(u_{1}, u_{2}, \cdots, u_{l}, 0, \cdots, 0) = \boldsymbol{f} \circ h(u_{1}, u_{2}, \cdots, u_{n}),
\]
\[
  \left.\frac{\partial(G_{1}, G_{2}, \cdots, G_{m})}{\partial(u_{1}, u_{2}, \cdots, u_{m})}\right|_{\boldsymbol{0}} = \det\begin{pmatrix}
    \boldsymbol{I}_{l} & \boldsymbol{0} \\
    * & \boldsymbol{I}_{m-l}
  \end{pmatrix}_{\boldsymbol{0}} = 1 \neq 0.
\]
根据\cref{thm:8.5.1}, $G$ 有局部逆射 $g$, 将 $\mathbb{R}^{m}$ 中的 $\boldsymbol{0}$ 的一个开邻域映成 $\boldsymbol{0}$ 的另一个开邻域. 并且
\[
  \boldsymbol{g} \circ \boldsymbol{f} \circ \boldsymbol{h}(x_{1}, x_{2}, \cdots, x_{n}) = \boldsymbol{g} \circ \boldsymbol{G}(x_{1}, x_{2}, \cdots, x_{l}, 0, \cdots, 0) = (x_{1}, x_{2}, \cdots, x_{l}, 0, \cdots, 0).
\]
\end{proof}
\begin{exblock}{复习题 8}
\item 设 $f_{x}^{\prime}$, $f_{y}^{\prime}$ 与 $f_{yx}^{\prime\prime}$ 在点 $(x_{0}, y_{0})$ 的某开邻域内存在, $f_{yx}^{\prime\prime}$ 在点 $(x_{0}, y_{0})$ 连续, 证明: $f_{xy}^{\prime\prime}(x_{0}, y_{0})$ 也存在, 且 $f_{xy}^{\prime\prime}(x_{0}, y_{0}) = f_{xy}^{\prime\prime}(x_{0}, y_{0})$.

\item 设 $f_{x}^{\prime}$, $f_{y}^{\prime}$ 在点 $(x_{0}, y_{0})$ 的某开邻域内存在且在点 $(x_{0}, y_{0})$ 可微, 证明:
\[
  f_{xy}^{\prime\prime}(x_{0}, y_{0}) = f_{yx}^{\prime\prime}(x_{0}, y_{0}).
\]

\item 设
\[
  u = \begin{vmatrix}
    1 & 1 & \cdots & 1 \\
    x_{1} & x_{2} & \cdots & x_{n} \\
    x_{1}^{2} & x_{2}^{2} & \cdots & x_{n}^{2} \\
    \vdots & \vdots & & \vdots \\
    x_{1}^{n-1} & x_{2}^{n-1} & \cdots & x_{n}^{n-1}
  \end{vmatrix},
\]
证明:
\begin{enumerate}
  \item[(1)] $\displaystyle\sum_{i=1}^{n}\frac{\partial u}{\partial x_{i}} = 0$;
  \item[(2)] $\displaystyle\sum_{i=1}^{n}x_{i}\frac{\partial u}{\partial x_{i}} = \frac{n(n-1)}{2}u$
\end{enumerate}

\item 设函数 $f(x,y)$ 具有连续的 $n$ 阶偏导数, 试证: 函数 $g(t) = f(a + ht, b + kt)$ 的 $n$ 阶导数
\[
  \frac{\mathrm{d}^{n}g(t)}{\mathrm{d}t^{n}} = \left(h\frac{\partial}{\partial x} + k\frac{\partial}{\partial y}\right)^{n}f(a + ht, b + kt).
\]

\item 设 $f(x,y)$ 为 $n$ 次齐次函数, 即对 $\forall\,t > 0$, 有 $f(tx, ty) = t^{n}f(x,y)$. 证明:
\[
  \left(x\frac{\partial}{\partial x} + y\frac{\partial}{\partial y}\right)^{m}f = n(n-1)\cdots(n-m+1)f.
\]

\item 对于函数 $f(x,y) = \sin\dfrac{y}{x}$, 证明: $\left(x\dfrac{\partial}{\partial x} + y\dfrac{\partial}{\partial y}\right)^{m}f = 0$.

\item 设 $\varphi(x,y)$ 为 $C^{2}$ 函数, 它满足弦振动方程: $a^{2}\dfrac{\partial^{2}\varphi}{\partial x^{2}} = \dfrac{\partial^{2}\varphi}{\partial t^{2}}$ $(a > 0)$.
\begin{enumerate}
  \item[(1)] 作变换 $u = x + at$, $v = x - at$, 则有 $\dfrac{\partial^{2}\varphi}{\partial u \partial v} = 0$;
  \item[(2)] 解得 $\varphi = f(u) + g(v) = f(x + at) + g(x - at)$
\end{enumerate}

\item 设 $u = u(x, y, z)$, $v = v(x, y, z)$ 和 $x = x(s, t)$, $y = y(s, t)$, $z = z(s, t)$ 都有连续的一阶偏导数. 证明:
\[
  \frac{\partial(u,v)}{\partial(s,t)} = \frac{\partial(u,v)}{\partial(x,y)}\frac{\partial(x,y)}{\partial(s,t)} + \frac{\partial(u,v)}{\partial(y,z)}\frac{\partial(y,z)}{\partial(s,t)} + \frac{\partial(u,v)}{\partial(z,x)}\frac{\partial(z,x)}{\partial(s,t)}.
\]

\item 设 $u = \dfrac{x}{r^{2}}$, $v = \dfrac{y}{r^{2}}$, $w = \dfrac{z}{r^{2}}$, 其中 $r = \sqrt{x^{2} + y^{2} + z^{2}}$. 证明:
\begin{enumerate}
  \item[(1)] 以 $u$, $v$, $w$ 为自变量的反函数组为
  \[
    x = \frac{u}{u^{2} + v^{2} + w^{2}}, \quad y = \frac{v}{u^{2} + v^{2} + w^{2}}, \quad z = \frac{w}{u^{2} + v^{2} + w^{2}};
  \]
  \item[(2)] $\displaystyle\dfrac{\partial(u,v,w)}{\partial(x,y,z)} = -\dfrac{1}{r^{6}}$
\end{enumerate}

\item 证明: Descartes (笛卡儿) 叶形线
\[
  x^{3} + y^{3} - 3axy = 0 \quad (a > 0)
\]
所确定的隐函数 $y = f(x)$ 的一阶与二阶导数分别为
\[
  y^{\prime} = \frac{ay - x^{2}}{y^{2} - ax} \quad (y^{2} - ax \neq 0)
\]
与
\[
  y^{\prime\prime} = -\frac{2a^{3}xy}{(y^{2} - ax)^{2}}.
\]
由此给出曲线上具有水平切线与垂直切线的点.

\item 求由方程 $z = f(x + y + z, xyz)$ 所确定的隐函数的偏导数 $\dfrac{\partial z}{\partial x}$, $\dfrac{\partial x}{\partial y}$, $\dfrac{\partial y}{\partial z}$, 并叙述可计算的条件.

\item 设 $u = x^{2} + y^{2} + z^{2}$, 其中 $z = f(x, y)$ 是由方程 $x^{3} + y^{3} + z^{3} = 3xyz$ 所确定的隐函数. 求 $u_{x}^{\prime}$ 与 $u_{xx}^{\prime\prime}$.

\item 应用球坐标与直角坐标之间的变换公式
\[
  (x,y,z) = \bigl(r\sin\theta\cos\varphi, r\sin\theta\sin\varphi, r\cos\theta\bigr),
\]
证明:
\[
  \Delta_{1}u = \left(\frac{\partial u}{\partial x}\right)^{2} + \left(\frac{\partial u}{\partial y}\right)^{2} + \left(\frac{\partial u}{\partial z}\right)^{2} = \left(\frac{\partial u}{\partial r}\right)^{2} + \frac{1}{r^{2}}\left(\frac{\partial u}{\partial \theta}\right)^{2} + \frac{1}{r^{2}\sin^{2}\theta}\left(\frac{\partial u}{\partial \varphi}\right)^{2};
\]
\[
  \Delta_{2}u = \frac{\partial^{2}u}{\partial x^{2}} + \frac{\partial^{2}u}{\partial y^{2}} + \frac{\partial^{2}u}{\partial z^{2}} = \frac{1}{r^{2}}\left[\frac{\partial}{\partial r}\left(r^{2}\frac{\partial u}{\partial r}\right) + \frac{1}{\sin\theta}\frac{\partial}{\partial\theta}\left(\sin\theta\frac{\partial u}{\partial\theta}\right) + \frac{1}{\sin^{2}\theta}\frac{\partial^{2}u}{\partial\varphi^{2}}\right].
\]
(提示: 上述变换是由两个特殊变换 $(x,y,z) = (R\cos\varphi, R\sin\varphi, z)$ 与 $(R,\varphi,z) = (r\sin\theta, \varphi, r\cos\theta)$ 复合而成)

\item 设 $f(x,y)$ 在点 $(x_{0}, y_{0})$ 的某开邻域内有偏导数 $f_{x}^{\prime}(x,y)$ 与 $f_{y}^{\prime}(x,y)$, 且 $f_{y}^{\prime}(x,y)$ 在点 $(x_{0}, y_{0})$ 处连续. 证明: $f(x,y)$ 在 $(x_{0}, y_{0})$ 处可微.

\item 设 $f(x,y) = |x-y|\varphi(x,y)$, 其中 $\varphi(x,y)$ 在点 $(0,0)$ 的某个开邻域内连续. 问:
\begin{enumerate}
  \item[(1)] $\varphi(x,y)$ 在什么条件下, 偏导数 $f_{x}^{\prime}(0,0)$, $f_{y}^{\prime}(0,0)$ 存在;
  \item[(2)] $\varphi(x,y)$ 在什么条件下, $f(x,y)$ 在 $(0,0)$ 处可微?
\end{enumerate}

\item 设函数 $f(x,y)$ 在闭单位圆片 $\bigl\{(x,y) \mid x^{2} + y^{2} \leqslant 1\bigr\}$ 上有连续的偏导数, 并且 $f(1,0) = f(0,1)$. 证明: 在单位圆 $\bigl\{(x,y) \mid x^{2} + y^{2} = 1\bigr\}$ 上至少有两点满足方程:
\[
  y\frac{\partial}{\partial x}f(x,y) = x\frac{\partial}{\partial y}f(x,y).
\]

\item 设实数 $x$, $y$, $z$ 满足 $e^{x} + y^{2} + |z| = 3$. 证明: $e^{x}y^{2}|z| \leqslant 1$.

\item 考察变换
\[
  \begin{cases}
    x = a_{1}u + b_{1}v + c_{1}w, \\
    y = a_{2}u + b_{2}v + c_{2}w, \\
    z = a_{3}u + b_{3}v + c_{3}w.
  \end{cases}
\]
问: 在什么条件下 (即 $a_{i}$, $b_{i}$, $c_{i}$ 满足什么条件时), 对任何二阶连续可微函数 $f$, $\left(\dfrac{\partial f}{\partial x}\right)^{2} + \left(\dfrac{\partial f}{\partial y}\right)^{2} + \left(\dfrac{\partial f}{\partial z}\right)^{2}$ 与 $\dfrac{\partial^{2}f}{\partial x^{2}} + \dfrac{\partial^{2}f}{\partial y^{2}} + \dfrac{\partial^{2}f}{\partial z^{2}}$ 在此变换下形式不变, 即
\[
  \left(\frac{\partial f}{\partial x}\right)^{2} + \left(\frac{\partial f}{\partial y}\right)^{2} + \left(\frac{\partial f}{\partial z}\right)^{2} = \left(\frac{\partial f}{\partial u}\right)^{2} + \left(\frac{\partial f}{\partial v}\right)^{2} + \left(\frac{\partial f}{\partial w}\right)^{2},
\]
\[
  \frac{\partial^{2}f}{\partial x^{2}} + \frac{\partial^{2}f}{\partial y^{2}} + \frac{\partial^{2}f}{\partial z^{2}} = \frac{\partial^{2}f}{\partial u^{2}} + \frac{\partial^{2}f}{\partial v^{2}} + \frac{\partial^{2}f}{\partial w^{2}}.
\]

\item 设 $f_{x}^{\prime}$, $f_{y}^{\prime}$ 在 $(x_{0}, y_{0})$ 的某开邻域内存在, 且在点 $(x_{0}, y_{0})$ 处可微. 证明:
\[
  f_{xy}^{\prime\prime}(x_{0}, y_{0}) = f_{yx}^{\prime\prime}(x_{0}, y_{0}).
\]

\item 设 $F(x,y,z)$ 在 $\mathbb{R}^{3}$ 中有连续的一阶偏导数 $\dfrac{\partial F}{\partial x}$, $\dfrac{\partial F}{\partial y}$, $\dfrac{\partial F}{\partial z}$, 并满足不等式:
\[
  y\frac{\partial F}{\partial x} - x\frac{\partial F}{\partial y} + \frac{\partial F}{\partial z} \geqslant \alpha > 0, \quad \forall (x,y,z) \in \mathbb{R}^{3},
\]
其中 $\alpha$ 为常数. 证明: 当 $(x,y,z)$ 沿着曲线
\[
  C: (x,y,z) = (-\cos t, \sin t, t), \quad t \geqslant 0
\]
趋向无穷远时, $F(x,y,z) \to +\infty$

\item 设 $U \subset \mathbb{R}^{n}$ 为凸区域, $f(\boldsymbol{x}) = f(x_{1}, x_{2}, \cdots, x_{n})$ 在 $U$ 上有连续的二阶偏导数. 证明: $f(\boldsymbol{x})$ 在 $U$ 上为凸函数等价于 $f$ 的 Hesse 矩阵
\[
  \boldsymbol{H}f = \left(\frac{\partial^{2}f}{\partial x_{i}\partial x_{j}}\right)
\]
在 $U$ 上是半正定的.
\end{exblock}
